\subsection{内正则集合函数}

定义 1. --- 设 T 是拓扑空间，\(\mathfrak{B}(T)\) 是 T 的 Borel σ-集环；令 I 是从 \(\mathfrak{B}(T)\) 到 \(\overline{\mathbf{R}}_{+}\) 的映射。

\begin{enumerate}
\def\labelenumi{\alph{enumi})}
\tightlist
\item
  称 I 是可数可加的，如果对于每个序列 \((\mathbf{A}_{n})\)（其元素是 \(\mathfrak{B}(\mathbf{T})\) 中两两不交的），
\end{enumerate}

\[
\mathrm{I} \left(\bigcup_ {n} \mathrm{A} _ {n}\right) = \sum_ {n} \mathrm{I} (\mathrm{A} _ {n}).\tag{4}
\]

\begin{enumerate}
\def\labelenumi{\alph{enumi})}
\setcounter{enumi}{1}
\tightlist
\item
  称 I 是内正则的，如果对于每个集合 \(\mathbf{A} \in \mathfrak{B}(\mathbf{T})\)，
\end{enumerate}

\[
\mathrm{I} (\mathrm{A}) = \sup _ {\mathrm{K}} \mathrm{I} (\mathrm{K}),\tag{5}
\]

其中 K 遍历 A 的紧子集所成的集合。

\begin{enumerate}
\def\labelenumi{\alph{enumi})}
\setcounter{enumi}{2}
\tightlist
\item
  称 I 有界（相应地，局部有界），如果 \(\mathrm{I(T)} < +\infty\)（相应地，如果 T 中每一点 \(x \in \mathbf{T}\) 都有一个开邻域 V，使得 \(\mathrm{I(V)} < +\infty\)）。
\end{enumerate}

注记。--- 1) 条件 a) 显然蕴含 I 是从 \(\mathfrak{B}(\mathbf{T})\)（按包含关系排序）到 \(\overline{\mathbf{R}}_{+}\) 的递增映射。

\begin{enumerate}
\def\labelenumi{\arabic{enumi})}
\setcounter{enumi}{1}
\item
  假设 I 可数可加；令 \((\mathbf{A}_n)_{n\in \mathbf{N}}\) 是 Borel 集的递增序列，并令 \(\mathbf{A} = \bigcup_{n\in \mathbf{N}}\mathbf{A}_n\)。由于集合 \(\mathrm{D}_0 = \mathrm{A}_0\)、\(\mathrm{D}_n = \mathrm{A}_n - \mathrm{A}_{n - 1}\) 两两不交且其并为 A，有 \(\mathrm{I}(\mathrm{A}) = \sum_{n} \mathrm{I}(\mathrm{D}_{n}) = \lim_{n \to \infty} \mathrm{I}(\mathrm{A}_{n})\)。同样，如果 \((\mathrm{B}_{n})\) 是 Borel 集的递减序列且 \(\mathrm{I}(\mathrm{B}_{0}) < +\infty\)，则 \(\mathrm{I}\left(\bigcap_{n} \mathrm{B}_{n}\right) = \lim_{n \to \infty} \mathrm{I}(\mathrm{B}_{n})\)：只需将前面的结论应用于集合 \(A_{n} = B_{0} - B_{n}\)。
\item
  令 \((\mathbf{A}_n)\) 是 T 的 Borel 集的任意序列。如果 I 可数可加，则 \(\mathrm{I}\left(\bigcup_{n} \mathrm{A}_n\right) \leqslant \sum_{n} \mathrm{I}(\mathrm{A}_n)\)。根据前一个注记，只需对有限序列证明这个不等式。立即可归结为两个集合 \(A_{1}\) 和 \(A_{2}\) 的情形；但关系 (4) 蕴含
\end{enumerate}

\[
\mathrm{I} \left(\mathrm{A} _ {1} \cup \mathrm{A} _ {2}\right) = \mathrm{I} \left(\mathrm{A} _ {1}\right) + \mathrm{I} \left(\mathrm{A} _ {2} - \left(\mathrm{A} _ {1} \cap \mathrm{A} _ {2}\right)\right) \leqslant \mathrm{I} \left(\mathrm{A} _ {1}\right) + \mathrm{I} \left(\mathrm{A} _ {2}\right).
\]

所需不等式立即得到。

\begin{enumerate}
\def\labelenumi{\arabic{enumi})}
\setcounter{enumi}{3}
\item
  如果 I 是可数可加且局部有界的函数，则前一个注记立即蕴含，有 \(\mathrm{I}(\mathbf{K}) < +\infty\)，对于每个紧集 \(\mathbf{K} \subset \mathbf{T}\)。
\item
  可以证明，如果 I 可加，即对有限序列满足 (4)，且 I 内正则，则 I 可数可加（Exer. 7）。读者还可验证，下面 Th. 2 的证明只用到了可加性和内正则性。
\end{enumerate}

定理 2. --- 设 T 是拓扑空间，I 是定义在 \(\mathfrak{B}(T)\) 上、取值于 \(\overline{\mathbf{R}}_{+}\) 的函数。要使存在 T 上的测度 \(\mu\)，使得 \(\mu^{\bullet}(\mathbf{A}) = \mathrm{I}(\mathbf{A})\) 对于每个 \(\mathbf{A} \in \mathfrak{B}(T)\) 成立，必要且充分的条件是 I 可数可加、局部有界且内正则。测度 \(\mu\) 于是唯一。

这三个条件是必要的：事实上，映射 \(A \mapsto \mu^{\bullet}(A)\) 定义在 \(\mathfrak{B}(T)\) 上，是可数可加的（\(\S1\) , No.~5, Cor. of Prop. 4），由测度的定义它是局部有界的（\(\S1\) , No.~2, Def. 5），并且根据 \(\S1\) , No.~2 的注记 3 它是内正则的。

下面证明存在性。显然，I 在 \(\mathcal{K}(\mathrm{T})\) 上的限制满足 Th. 1 叙述中的条件 1)、2)、3) 和 5)；我们还要证明它满足 4)。设 K 是 T 的紧子集，是紧集递减有向族 \((\mathrm{K}_{\alpha})_{\alpha \in \mathrm{A}}\) 的交，并令 \(\varepsilon\) \(>0\)；由于 I 局部有界，存在 K 的一个开（因而是 Borel）邻域 V，使得 \(\mathrm{I(V)} < +\infty\)，于是存在指标 \(\alpha\)，使得 \(\mathrm{K}_{\alpha} \subset \mathrm{V}\)；必要时改变记号，可以假设 \(\mathrm{K}_{\alpha} \subset \mathrm{V}\) 对所有 \(\alpha \in \mathrm{A}\) 成立。由 I 的内正则性，存在紧集 \(\mathrm{L} \subset \mathrm{V} - \mathrm{K}\)，使得 \(\mathrm{I(L)} \geqslant \mathrm{I(V} - \mathrm{K)} - \varepsilon\)；由于 L 不与 K 相交，存在指标 \(\alpha\)，使得 \(\mathrm{L} \cap \mathrm{K}_{\alpha} = \varnothing\)，于是有 \(\mathrm{I}(\mathrm{V} - \mathrm{K}_{\alpha}) \geqslant \mathrm{I}(\mathrm{L}) \geqslant \mathrm{I}(\mathrm{V} - \mathrm{K}) - \varepsilon\)。由于 \(\mathbf{K}_{\alpha} \subset \mathbf{V}\)，可得 \(\mathrm{I}(\mathbf{K}_{\alpha}) \leqslant \mathrm{I}(\mathbf{K}) + \varepsilon\)，条件 4) 得到验证。

根据 Th. 1，存在测度 \(\mu\)，使得 \(\mu^{\bullet}(\mathrm{K}) = \mathrm{I}(\mathrm{K})\) 对所有 \(\mathbf{K} \in \mathfrak{K}(\mathbf{T})\) 成立。集合函数 \(\mu^{\bullet}\) 和 I 在 \(\mathfrak{B}(\mathbf{T})\) 上的内正则性于是蕴含 \(\mu^{\bullet}(\mathbf{A}) = \mathrm{I}(\mathbf{A})\) 对所有 \(\mathbf{A} \in \mathfrak{B}(\mathbf{T})\) 成立，存在性得证。\(\mu\) 的唯一性由 Th. 1 的唯一性断言得出。

